\begin{theorem}[IntUp signed-history RatUp numerator scope]
\label{thm:int-history-certificate-rat-sibling-scope}
Every $\RatUp$ numerator readback that consumes an $\IntUp$ signed-history row factors through the rational naming certificate of \autoref{ch:concrete-instances-rat-namecert}. The consumed integer row is still the signed-unary package generated from the $\NatUp$ magnitude surface of \autoref{ch:concrete-instances-nat-namecert} and the continuation-compatible additive routing of \autoref{ch:concrete-instances-add-namecert}; the rational consumer merely places that displayed signed row beside its positive denominator and ledger rows.

Consequently the downstream rational read exports no integer quotient primitive, no field equality, and no rational equality row at the $\IntUp$ boundary. The only exported integer data are the visible sign mark, unary magnitude, zero-sign normalization, append-context transport, and signed-history classifier rows already supplied by the $\IntUp$ certificate.
\end{theorem}
\begin{proof}
Read the signed-history certificate through \autoref{thm:intup-mature-signed-unary-downstream-package}. That package exposes the semantic certificate, signed-history classifier rows, append-context and $\hsame$ transports, zero-sign boundary, and standard signed-unary bridge.

The rational consumer is then the carrier of \autoref{ch:concrete-instances-rat-namecert}. It accepts the displayed signed integer row as numerator data and pairs it with the positive denominator and rational ledger rows, so the readback factors through the rational sibling chapter rather than through an ambient quotient construction.

The source rows for the integer side remain the Nat unary magnitude row and the Add continuation route. Since no extra coordinate is introduced at the handoff, the boundary exports neither an integer quotient primitive nor a field-level equality witness.
\end{proof}
